\begin{frame}{0. Architecture C++ / Interface Python et ouverture du
code}
\protect\hypertarget{architecture-c-interface-python-et-ouverture-du-code}{}
\begin{block}{\textbf{Base C++}}
\protect\hypertarget{base-c}{}
\begin{itemize}
\tightlist
\item
  Performances élevées pour les opérations lourdes sur les maillages.
\item
  Intégration naturelle avec les codes de calcul existants.
\item
  Gestion fine de la mémoire, adaptée au HPC.
\end{itemize}
\end{block}

\begin{block}{\textbf{Interface Python}}
\protect\hypertarget{interface-python}{}
\begin{itemize}
\tightlist
\item
  Utilisation quotidienne simple pour les études, le pré/post-traitement
  et le prototypage.
\item
  Accès direct aux structures C++ sans perte de performance sur les
  calculs.
\item
  Intégration facile dans les workflows Python déjà largement utilisés.
\end{itemize}
\end{block}

\begin{block}{\textbf{Ouverture du code}}
\protect\hypertarget{ouverture-du-code}{}
\begin{itemize}
\tightlist
\item
  Code open source, facile à intégrer dans des projets variés du CEA ou
  en collaboration externe (EdF).
\item
  Transparence sur les algorithmes, possibilité d'amélioration
  collaborative.
\item
  Favorise la pérennité et la mutualisation entre équipes.
\end{itemize}
\end{block}

\begin{block}{\textbf{Intérêt dans les projets}}
\protect\hypertarget{intuxe9ruxeat-dans-les-projets}{}
\begin{itemize}
\tightlist
\item
  À destination de développeurs de pré-pro complexes construits en
  python (comethe pour Flica5, Apollo)
\item
  À destination de développeurs de codes de calcul pour des besoins
  avancés (TRUST, Apollo pour le couplage)
\item
  il ne s'agit pas d'un outil disponible avec une interface graphique
  pour un utilisateur qui réalise des études.
\end{itemize}
\end{block}
\end{frame}

\begin{frame}{1. Génération procédurale et manipulation de maillages}
\protect\hypertarget{guxe9nuxe9ration-procuxe9durale-et-manipulation-de-maillages}{}
MEDCoupling fournit un ensemble d'outils permettant de créer,
transformer et analyser des maillages. Ces fonctionnalités couvrent les
besoins classiques des outils de pré/post-traitement pour plusieurs
physiques (thermo-hydraulique, mécanique, neutronique).

\begin{block}{Principales capacités}
\protect\hypertarget{principales-capacituxe9s}{}
\begin{block}{\textbf{Construction automatique de maillages}}
\protect\hypertarget{construction-automatique-de-maillages}{}
\begin{itemize}
\tightlist
\item
  maillages cartésiens et extrudés (1D, 2D, 3D) ;
\item
  maillages non structurés (triangles, quadrangles, tétraèdres,
  hexaèdres, polyèdres)
\item
  génération programmatique à partir de lignes, contours, surfaces ou
  volumes.
\end{itemize}
\end{block}

\begin{block}{\textbf{Opérations de modification}}
\protect\hypertarget{opuxe9rations-de-modification}{}
\begin{itemize}
\tightlist
\item
  concaténation, aggrégation, conformisation
\item
  transformations géométriques (translation, rotation, déformation).
\end{itemize}
\end{block}

\begin{block}{\textbf{Opérations géométriques avancées}}
\protect\hypertarget{opuxe9rations-guxe9omuxe9triques-avancuxe9es}{}
\begin{itemize}
\tightlist
\item
  calcul d'intersections (2D/2D, 3D/plan, etc) entre maillages ;
\item
  localisation de points ;
\item
  calcul de volumes, surfaces, centres de gravité.
\end{itemize}
\end{block}

\begin{block}{\textbf{Gestion des champs}}
\protect\hypertarget{gestion-des-champs}{}
\begin{itemize}
\tightlist
\item
  définition de champs nodaux (P1) ou cellulaires (P0) ;
\item
  support des champs intensifs/extensifs ;
\item
  prise en charge du temps et de l'évolution transitoire.
\end{itemize}
\end{block}

\begin{block}{\textbf{Export direct en VTK (format standard de
maillage)}}
\protect\hypertarget{export-direct-en-vtk-format-standard-de-maillage}{}
\begin{itemize}
\tightlist
\item
  pour la visualisation ou le post-traitement.
\end{itemize}
\end{block}
\end{block}

\begin{block}{Intérêt pour les projets}
\protect\hypertarget{intuxe9ruxeat-pour-les-projets}{}
\begin{itemize}
\tightlist
\item
  Permet de remplacer des scripts hétérogènes par une API python
  unifiée.
\item
  Propose des algorithmes efficaces en C++.
\item
  Offre une base commune entre équipes simulation, CFD, mécanique,
  neutronique, etc.
\end{itemize}
\end{block}
\end{frame}

\begin{frame}{2. Interface avec la librairie MED (EDF) pour la
lecture/écriture}
\protect\hypertarget{interface-avec-la-librairie-med-edf-pour-la-lectureuxe9criture}{}
MEDLoader est un module qui repose sur la bibliothèque MED développée
par EDF. Son rôle est de fournir une couche facilitant la
\textbf{conversion des données MED vers les objets internes de
MEDCoupling} et réciproquement. Il s'agit d'un outil de traduction entre
deux représentations différentes : la représentation du format MED et la
représentation interne simplifiée utilisée par MEDCoupling.

\begin{block}{Fonctionnalités principales}
\protect\hypertarget{fonctionnalituxe9s-principales}{}
MEDLoader expose une interface permettant :

\begin{itemize}
\tightlist
\item
  la lecture et l'écriture de maillages, champs, groupes, familles et
  profils,
\item
  la gestion du multi-pas de temps tel que défini dans le format MED,
\item
  l'écriture incrémentale (append),
\item
  la prise en charge des maillages multi-niveaux (maillage 3D, maillage
  de faces, maillage d'arêtes partageant les mêmes coordonnées),
\end{itemize}

Ces opérations reposent sur la conversion des données issues de MED vers
les structures MEDCoupling plus simples (DataArray, Mesh, Field), ce qui
facilite ensuite les manipulations géométriques ou les projections.
\end{block}

\begin{block}{Concepts manipulés}
\protect\hypertarget{concepts-manipuluxe9s}{}
\begin{itemize}
\tightlist
\item
  \textbf{Groupes} : ensembles nommés de cellules ou entités, souvent
  utilisés pour les conditions limites.
\item
  \textbf{Familles} : mécanisme interne du format MED permettant
  d'encoder la composition des groupes.
\item
  \textbf{Profils} : sous-ensembles d'indices, permettant de restreindre
  un champ à une partie du maillage.
\end{itemize}
\end{block}

\begin{block}{Intérêt pour les projets}
\protect\hypertarget{intuxe9ruxeat-pour-les-projets-1}{}
\begin{itemize}
\tightlist
\item
  Permet de réutiliser les structures MEDCoupling dans des workflows où
  le stockage se fait en format MED.
\item
  Facilite la mise en place de chaînes de traitement impliquant
  MEDCoupling (manipulations géométriques, projections,
  transformations).
\item
  Garantit la compatibilité avec les outils de visualisation habituels,
  notamment via l'export VTK.
\item
  Sert d'intermédiaire lorsqu'un code ou script doit convertir des
  données MED vers un format directement exploitable dans MEDCoupling.
\end{itemize}
\end{block}
\end{frame}

\begin{frame}{3. Couplage de codes (séquentiel et parallèle)}
\protect\hypertarget{couplage-de-codes-suxe9quentiel-et-paralluxe8le}{}
MEDCoupling offre des outils dédiés au \textbf{transfert de champs entre
codes}, que ce soit en mode séquentiel ou via MPI en parallèle.

Ces outils sont essentiels dans les couplages multiphysiques
(neutronique - thermique, fluide - thermique, fluide - structure, etc.),
en particulier lorsque les codes utilisent des maillages différents.

\begin{block}{3.1 Couplage séquentiel (Remapper)}
\protect\hypertarget{couplage-suxe9quentiel-remapper}{}
Le \textbf{Remapper} assure la projection de champs entre maillages,
même de nature très différente.

\begin{block}{Fonctionnement général}
\protect\hypertarget{fonctionnement-guxe9nuxe9ral}{}
\begin{enumerate}
\tightlist
\item
  \textbf{Préparation} : construction de la correspondance géométrique
  entre maillage source et maillage cible.
\item
  \textbf{Transfert} : application d'une projection conservatrice ou
  non, selon la nature du champ.
\end{enumerate}
\end{block}

\begin{block}{Choix de la nature du champ}
\protect\hypertarget{choix-de-la-nature-du-champ}{}
\begin{itemize}
\tightlist
\item
  \textbf{Intensive} : indépendantes de la maille (température,
  vitesse).
\item
  \textbf{Extensive} : proportionnelles au volume/aire (énergie, masse).
\item
  Modes : \emph{conservation}, \emph{maximum}, \emph{moyenne} selon les
  besoins physiques.
\end{itemize}
\end{block}

\begin{block}{Compatibilité}
\protect\hypertarget{compatibilituxe9}{}
\begin{itemize}
\tightlist
\item
  Maillages structurés, non structurés ou extrudés.
\item
  Dimensions 1D, 2D, 3D et surfaces 2D dans un espace 3D.
\item
  Champs P0, P1 ou points de Gauss.
\end{itemize}
\end{block}
\end{block}

\begin{block}{3.2 Couplage parallèle (DEC / ParaMEDMEM)}
\protect\hypertarget{couplage-paralluxe8le-dec-paramedmem}{}
Pour les applications HPC, MEDCoupling gère le transfert de champs entre
codes parallèles via \textbf{Data Exchange Channels (DEC)}.

\begin{block}{Principes de fonctionnement}
\protect\hypertarget{principes-de-fonctionnement}{}
\begin{itemize}
\tightlist
\item
  Échange des boîtes englobantes entre processus MPI.
\item
  Détermination des sous-domaines qui se recoupent.
\item
  Calcul des projections localement, sans reconstituer le maillage
  global.
\item
  Construction distribuée de la matrice de transfert.
\item
  Communication optimisée pour minimiser les échanges inter-processus.
\end{itemize}
\end{block}

\begin{block}{Intérêt pour les projets}
\protect\hypertarget{intuxe9ruxeat-pour-les-projets-2}{}
\begin{itemize}
\tightlist
\item
  Permet le couplage performant de codes massivement parallèles.
\item
  Évite la duplication des maillages complets sur chaque processus.
\item
  Maintient la cohérence du transfert même sur des domaines fortement
  découpés.
\end{itemize}
\end{block}
\end{block}
\end{frame}

\begin{frame}{Conclusion}
\protect\hypertarget{conclusion}{}
\begin{itemize}
\tightlist
\item
  Outil central pour plusieurs projets de simulation au CEA.
\item
  Noyau C++ performant et interface Python favorable au développement
  rapide.
\item
  Approche open source facilitant l'intégration et la collaboration.
\item
  Large éventail de fonctionnalités : génération, manipulation et
  analyse de maillages.
\item
  Socle commun pour les chaînes de pré/post-traitement.
\item
  Outils indispensable pour le transfert de champs et le couplage
  multiphysique, en séquentiel comme en parallèle.
\end{itemize}
\end{frame}
